% 真实考试场景示例 - 高中数学月考试卷
% 包含：密封线、草稿纸、完整答案、评分标准
% 编译：xelatex 03-real-exam-scenario.tex

\documentclass{exam-zh}

% ===== 显示答案（教师版）=====
% 学生版：注释掉下面的 \examsetup
% 教师版：取消注释下面的 \examsetup
% \examsetup{
%   question/show-answer = true,
%   solution/show-solution = show-stay,
%   question/show-points = true
% }

\begin{document}

% ===== 试卷密封线 =====
\secret  % 生成密封线

% ===== 试卷抬头 =====
\title{高一数学月考试卷}
\subject{数学}
\maketitle

\information{
  姓名\underline{\hspace{6em}},
  班级\underline{\hspace{6em}},
  学号\underline{\hspace{6em}},
  考场\underline{\hspace{4em}},
  座号\underline{\hspace{3em}}
}

\begin{notice}
  \item 本试卷共 4 页，22 小题，满分 150 分。考试用时 120 分钟。
  \item 答卷前，考生务必将自己的姓名、学号、班级填写在答题卡上。
  \item 作答选择题时，选出每小题答案后，用 2B 铅笔在答题卡上对应题目选项的答案信息点涂黑；如需改动，用橡皮擦干净后，再选涂其他答案。
  \item 非选择题必须用黑色字迹的钢笔或签字笔作答，答案必须写在答题卡各题目指定区域内相应位置上；如需改动，先划掉原来的答案，然后再写上新答案。
  \item 考生必须保持答题卡的整洁。考试结束后，将试卷和答题卡一并交回。
\end{notice}

\vspace{1em}

% ===== 一、选择题 =====
\section{选择题（本题共 12 小题，每小题 5 分，共 60 分。在每小题给出的四个选项中，只有一项是符合题目要求的）}

\begin{question}
  已知集合 $A = \{x \mid x^2 - 3x + 2 = 0\}$，$B = \{1, 2, 3\}$，则 $A \cap B = $ \paren[D]
  \begin{choices}
    \item $\{1\}$
    \item $\{2\}$
    \item $\{3\}$
    \item $\{1, 2\}$
  \end{choices}
\end{question}

\begin{question}
  函数 $f(x) = \ln(2x - 1)$ 的定义域为\paren[C]
  \begin{choices}
    \item $(-\infty, \frac{1}{2})$
    \item $[\frac{1}{2}, +\infty)$
    \item $(\frac{1}{2}, +\infty)$
    \item $\mathbb{R}$
  \end{choices}
\end{question}

\begin{question}
  若复数 $z$ 满足 $(1 + \iu)z = 2$，则 $|z| = $ \paren[B]
  \begin{choices}
    \item $1$
    \item $\sqrt{2}$
    \item $2$
    \item $2\sqrt{2}$
  \end{choices}
\end{question}

\begin{question}
  已知向量 $\vec{a} = (1, 2)$，$\vec{b} = (2, -1)$，则 $\vec{a} \cdot \vec{b} = $ \paren[A]
  \begin{choices}
    \item $0$
    \item $1$
    \item $-1$
    \item $3$
  \end{choices}
\end{question}

\begin{question}
  在等差数列 $\{a_n\}$ 中，$a_1 = 2$，$d = 3$，则 $a_5 = $ \paren[D]
  \begin{choices}
    \item $11$
    \item $13$
    \item $15$
    \item $14$
  \end{choices}
\end{question}

\begin{question}
  已知 $\sin \alpha = \frac{3}{5}$，且 $\alpha \in (\frac{\uppi}{2}, \uppi)$，则 $\cos \alpha = $ \paren[B]
  \begin{choices}
    \item $\frac{4}{5}$
    \item $-\frac{4}{5}$
    \item $\frac{3}{4}$
    \item $-\frac{3}{4}$
  \end{choices}
\end{question}

\begin{question}
  函数 $f(x) = 2^x - x^2$ 的零点所在的大致区间是\paren[C]
  \begin{choices}
    \item $(-1, 0)$
    \item $(0, 1)$
    \item $(1, 2)$
    \item $(2, 3)$
  \end{choices}
\end{question}

\begin{question}
  已知圆 $C: (x - 2)^2 + (y + 1)^2 = 5$，则圆心到直线 $3x + 4y - 5 = 0$ 的距离为\paren[A]
  \begin{choices}
    \item $1$
    \item $2$
    \item $\sqrt{5}$
    \item $5$
  \end{choices}
\end{question}

\begin{question}
  若 $\log_2 x = 3$，则 $x = $ \paren[C]
  \begin{choices}
    \item $3$
    \item $6$
    \item $8$
    \item $9$
  \end{choices}
\end{question}

\begin{question}
  在 $\triangle ABC$ 中，$a = 3$，$b = 4$，$C = 60^\circ$，则 $c = $ \paren[B]
  \begin{choices}
    \item $5$
    \item $\sqrt{13}$
    \item $\sqrt{37}$
    \item $7$
  \end{choices}
\end{question}

\begin{question}
  已知函数 $f(x) = x^3 - 3x$，则 $f(x)$ 在区间 $[-2, 2]$ 上的最小值为\paren[D]
  \begin{choices}
    \item $-4$
    \item $-3$
    \item $-1$
    \item $-2$
  \end{choices}
\end{question}

\begin{question}
  双曲线 $\frac{x^2}{4} - \frac{y^2}{5} = 1$ 的渐近线方程为\paren[A]
  \begin{choices}
    \item $y = \pm \frac{\sqrt{5}}{2}x$
    \item $y = \pm \frac{2}{\sqrt{5}}x$
    \item $y = \pm \frac{5}{4}x$
    \item $y = \pm \frac{4}{5}x$
  \end{choices}
\end{question}


% ===== 二、填空题 =====
\section{填空题（本题共 4 小题，每小题 5 分，共 20 分）}
\examsetup{question/index = 0}

\begin{question}
  已知 $f(x) = 2x^2 - 3x + 1$，则 $f(2) = $ \fillin[3]。
\end{question}

\begin{question}
  若 $\tan \alpha = 2$，则 $\frac{\sin \alpha + \cos \alpha}{\sin \alpha - \cos \alpha} = $ \fillin[3]。
\end{question}

\begin{question}
  已知等比数列 $\{a_n\}$ 满足 $a_1 = 1$，$a_3 = 4$，则公比 $q = $ \fillin[{$\pm 2$}]。
\end{question}

\begin{question}
  若不等式 $|x - 1| + |x + 2| \geq a$ 对所有实数 $x$ 恒成立，则实数 $a$ 的取值范围是 \fillin[{$(-\infty, 3]$}]。
\end{question}


% ===== 三、解答题 =====
\section{解答题（本题共 6 小题，共 70 分。解答应写出文字说明、证明过程或演算步骤）}
\examsetup{question/index = 0}

\begin{problem}[points = 10]
  已知函数 $f(x) = 2\sin x \cos x + 2\sqrt{3}\cos^2 x - \sqrt{3}$。

  \begin{enumerate}
    \item 求 $f(x)$ 的最小正周期；
    \item 求 $f(x)$ 的单调递增区间。
  \end{enumerate}

  \begin{solution}
    \begin{enumerate}
      \item $f(x) = \sin 2x + \sqrt{3}(2\cos^2 x - 1)$ \score{2}

      $= \sin 2x + \sqrt{3}\cos 2x$ \score{1}

      $= 2\sin(2x + \frac{\uppi}{3})$。\score{1}

      所以最小正周期 $T = \frac{2\uppi}{2} = \uppi$。\score{1}

      \item 由 $-\frac{\uppi}{2} + 2k\uppi \leq 2x + \frac{\uppi}{3} \leq \frac{\uppi}{2} + 2k\uppi$，$k \in \mathbb{Z}$，\score{2}

      得 $-\frac{5\uppi}{12} + k\uppi \leq x \leq \frac{\uppi}{12} + k\uppi$，$k \in \mathbb{Z}$。\score{2}

      所以 $f(x)$ 的单调递增区间为 $[-\frac{5\uppi}{12} + k\uppi, \frac{\uppi}{12} + k\uppi]$，$k \in \mathbb{Z}$。\score{1}
    \end{enumerate}
  \end{solution}
\end{problem}

\begin{problem}[points = 12]
  在 $\triangle ABC$ 中，角 $A$，$B$，$C$ 的对边分别为 $a$，$b$，$c$，已知 $b = 3$，$c = 2\sqrt{3}$，$B = 60^\circ$。

  \begin{enumerate}
    \item 求边 $a$ 的长；
    \item 求 $\triangle ABC$ 的面积。
  \end{enumerate}

  \begin{solution}
    \begin{enumerate}
      \item 由余弦定理 $b^2 = a^2 + c^2 - 2ac\cos B$，\score{2}

      得 $9 = a^2 + 12 - 2a \times 2\sqrt{3} \times \frac{1}{2}$，\score{2}

      即 $a^2 - 2\sqrt{3}a + 3 = 0$，\score{1}

      解得 $a = \sqrt{3}$ 或 $a = \sqrt{3}$（重根）。\score{1}

      所以 $a = \sqrt{3}$。\score{1}

      \item $S_{\triangle ABC} = \frac{1}{2}ac\sin B$ \score{2}

      $= \frac{1}{2} \times \sqrt{3} \times 2\sqrt{3} \times \frac{\sqrt{3}}{2}$ \score{2}

      $= \frac{3\sqrt{3}}{2}$。\score{1}
    \end{enumerate}
  \end{solution}
\end{problem}

\begin{problem}[points = 12]
  已知等差数列 $\{a_n\}$ 的前 $n$ 项和为 $S_n$，$a_2 = 4$，$S_5 = 30$。

  \begin{enumerate}
    \item 求数列 $\{a_n\}$ 的通项公式；
    \item 求满足 $S_n > 100$ 的最小正整数 $n$。
  \end{enumerate}

  \begin{solution}
    \begin{enumerate}
      \item 设首项为 $a_1$，公差为 $d$，由题意得\score{1}

      $\begin{cases} a_1 + d = 4 \\ 5a_1 + 10d = 30 \end{cases}$，\score{2}

      解得 $a_1 = 2$，$d = 2$。\score{2}

      所以 $a_n = 2 + 2(n - 1) = 2n$。\score{1}

      \item $S_n = \frac{n(a_1 + a_n)}{2} = \frac{n(2 + 2n)}{2} = n(n + 1)$。\score{2}

      由 $S_n > 100$ 得 $n(n + 1) > 100$，\score{1}

      解得 $n > \frac{-1 + \sqrt{401}}{2} \approx 9.51$。\score{2}

      所以最小正整数 $n = 10$。\score{1}
    \end{enumerate}
  \end{solution}
\end{problem}

\begin{problem}[points = 12]
  已知函数 $f(x) = x^3 - 3x^2 + 1$。

  \begin{enumerate}
    \item 求 $f(x)$ 的单调区间；
    \item 求 $f(x)$ 在区间 $[-1, 3]$ 上的最大值和最小值。
  \end{enumerate}

  \begin{solution}
    \begin{enumerate}
      \item $f'(x) = 3x^2 - 6x = 3x(x - 2)$。\score{2}

      令 $f'(x) = 0$，得 $x = 0$ 或 $x = 2$。\score{1}

      当 $x < 0$ 或 $x > 2$ 时，$f'(x) > 0$；\score{1}

      当 $0 < x < 2$ 时，$f'(x) < 0$。\score{1}

      所以 $f(x)$ 的递增区间为 $(-\infty, 0)$ 和 $(2, +\infty)$，

      递减区间为 $(0, 2)$。\score{1}

      \item $f(-1) = -3$，$f(0) = 1$，$f(2) = -3$，$f(3) = 1$。\score{3}

      所以 $f(x)$ 在 $[-1, 3]$ 上的最大值为 $1$，最小值为 $-3$。\score{3}
    \end{enumerate}
  \end{solution}
\end{problem}

\begin{problem}[points = 12]
  已知圆 $C: x^2 + y^2 - 2x - 4y + 1 = 0$。

  \begin{enumerate}
    \item 求圆 $C$ 的圆心坐标和半径；
    \item 若直线 $l: x + y + m = 0$ 与圆 $C$ 相切，求 $m$ 的值。
  \end{enumerate}

  \begin{solution}
    \begin{enumerate}
      \item 圆的方程可化为 $(x - 1)^2 + (y - 2)^2 = 4$。\score{3}

      所以圆心为 $(1, 2)$，半径 $r = 2$。\score{2}

      \item 圆心到直线的距离 $d = \frac{|1 + 2 + m|}{\sqrt{2}} = \frac{|3 + m|}{\sqrt{2}}$。\score{3}

      由相切条件 $d = r$，得 $\frac{|3 + m|}{\sqrt{2}} = 2$，\score{2}

      解得 $m = 2\sqrt{2} - 3$ 或 $m = -2\sqrt{2} - 3$。\score{2}
    \end{enumerate}
  \end{solution}
\end{problem}

\begin{problem}[points = 12]
  已知函数 $f(x) = \ln x - ax$（$a \in \mathbb{R}$）。

  \begin{enumerate}
    \item 讨论 $f(x)$ 的单调性；
    \item 当 $a = 1$ 时，证明：$f(x) \leq -1$。
  \end{enumerate}

  \begin{solution}
    \begin{enumerate}
      \item $f'(x) = \frac{1}{x} - a = \frac{1 - ax}{x}$（$x > 0$）。\score{2}

      当 $a \leq 0$ 时，$f'(x) > 0$ 恒成立，$f(x)$ 在 $(0, +\infty)$ 上单调递增。\score{2}

      当 $a > 0$ 时，令 $f'(x) = 0$，得 $x = \frac{1}{a}$。\score{1}

      当 $0 < x < \frac{1}{a}$ 时，$f'(x) > 0$；当 $x > \frac{1}{a}$ 时，$f'(x) < 0$。\score{1}

      所以 $f(x)$ 在 $(0, \frac{1}{a})$ 上单调递增，在 $(\frac{1}{a}, +\infty)$ 上单调递减。\score{1}

      \item 当 $a = 1$ 时，$f(x) = \ln x - x$。\score{1}

      由（1）知，$f(x)$ 在 $x = 1$ 处取得最大值 $f(1) = -1$。\score{2}

      所以 $f(x) \leq -1$。\score{2}
    \end{enumerate}
  \end{solution}
\end{problem}


% ===== 草稿纸 =====
\newpage
\section*{草稿纸}

\draftpaper

\end{document}
